\customchapter{TRABAJOS FUTUROS}

\begin{itemize}
\item \textbf{Repetición completa.} Entrenar HyperNCD con varias semillas y
repetir después la extracción y el agrupamiento, para separar la variabilidad
del aprendizaje de la variabilidad de la etapa final.

\item \textbf{Validación sin reutilizar el conjunto de prueba.} Reservar escenas
de ArCH para validación o aplicar validación cruzada, manteniendo intactas las
escenas oficiales de prueba para la evaluación final.

\item \textbf{Más escenas patrimoniales.} Evaluar iglesias, conventos, huacas y
edificios republicanos para determinar si la ventaja de ArCH se conserva fuera
de la sección estudiada de la Catedral de Lima.

\item \textbf{Ablación de la fusión.} Comparar bajo las mismas semillas el
descriptor semántico, Sonata, la geometría manual y sus combinaciones, e
incorporar intervalos de confianza o pruebas pareadas antes de atribuir la
diferencia a una rama específica.

\item \textbf{Clases minoritarias y ornamento fino.} Ajustar el muestreo, la
resolución y las pérdidas para puerta/ventana, arco y escaleras, que conservaron
los IoU más bajos durante el entrenamiento con ArCH.

\item \textbf{Vocabulario abierto e instancias.} Nombrar las agrupaciones
descubiertas y separar elementos individuales, como columna~1 y columna~2, para
facilitar su incorporación posterior a un modelo BIM.
\end{itemize}
